\documentclass[11pt,a4paper]{article}

% ── kodowanie i język ─────────────────────────────────────────────────────────
\usepackage[T1]{fontenc}
\usepackage[utf8]{inputenc}
\usepackage{lmodern}
\usepackage[polish]{babel}
\usepackage[margin=2.5cm]{geometry}
\usepackage{microtype}
\usepackage{setspace}
\usepackage{csquotes}

% ── grafika ───────────────────────────────────────────────────────────────────
\usepackage[dvipsnames,svgnames]{xcolor}
\usepackage{graphicx}
\usepackage{tikz}
\usetikzlibrary{arrows.meta, positioning, fit, calc, shapes.geometric, backgrounds}

% ── tabele i listy ────────────────────────────────────────────────────────────
\usepackage{booktabs}
\usepackage{tabularx}
\usepackage{enumitem}
\usepackage{amsmath, amssymb}

% ── odnośniki (ładowane na końcu) ─────────────────────────────────────────────
\usepackage{hyperref}

% ── stonowana paleta ──────────────────────────────────────────────────────────
\definecolor{Navy}{HTML}{1F3A5F}      % nagłówki
\definecolor{Blue}{HTML}{2F6DB5}      % akcent 1 (osoby / klasyczny RAG)
\definecolor{Teal}{HTML}{2E8B82}      % akcent 2 (organizacje / graf)
\definecolor{Amber}{HTML}{B8860B}     % zdarzenia
\definecolor{Slate}{HTML}{5A6472}     % lokalizacje / tekst pomocniczy
\definecolor{LightBlue}{HTML}{E4EDF7}
\definecolor{LightTeal}{HTML}{E0F0EE}
\definecolor{LightAmber}{HTML}{F6EFD9}
\definecolor{LightSlate}{HTML}{ECEEF1}

% odnośniki w stonowanym granacie, żeby referencje do rysunków były widoczne
\hypersetup{colorlinks=true, linkcolor=Navy, citecolor=Navy, urlcolor=Navy}

% ── nagłówki sekcji: jeden stonowany granat ───────────────────────────────────
\usepackage{titlesec}
\titleformat{\section}{\Large\bfseries\color{Navy}}{\thesection}{0.6em}{}
\titleformat{\subsection}{\large\bfseries\color{Navy}}{\thesubsection}{0.6em}{}
\titleformat{\subsubsection}{\normalsize\bfseries\color{Navy}}{\thesubsubsection}{0.6em}{}

\setlength{\parskip}{0.5em}
\setlength{\parindent}{0pt}

\begin{document}

% ── strona tytułowa ───────────────────────────────────────────────────────────
\begin{titlepage}
\centering
\vspace*{3cm}

{\color{Slate}\large Raport projektowy}\\[2.5cm]

{\color{Navy}\LARGE\bfseries Opracowanie prototypu GraphRAG\\[0.5em]
wspomagającego wyszukiwanie powiązań\\[0.5em]
pomiędzy obiektami i zdarzeniami\\[0.5em]
występującymi w artykułach prasowych}\\[3.5cm]

{\large
Miłosz Słowiński\\[0.4em]
Piotr Karamon
}\\[3cm]

{\color{Slate} 25 czerwca 2026}

\vfill
\end{titlepage}

\tableofcontents
\newpage

% ══════════════════════════════════════════════════════════════════════════════
\section{Klasyczny RAG i jego ograniczenia}
% ══════════════════════════════════════════════════════════════════════════════

\subsection{Jak działa zwykły RAG}

RAG (\textit{retrieval-augmented generation}, czyli generowanie wspomagane wyszukiwaniem)
to technika, w której model językowy dostaje dodatkową wiedzę z zewnętrznego źródła. Sam
model zna tylko to, czego nauczył się podczas treningu, i ma ograniczone okno kontekstu, więc
nie da się wkleić do niego całej bazy dokumentów. Zamiast tego do modelu trafiają tylko
niewielkie, dobrze dobrane fragmenty tekstu, które powinny zawierać odpowiedź.

W najpopularniejszym wariancie, nazywanym tu RAG wektorowym, działa to następująco:

\begin{enumerate}[itemsep=2pt]
  \item Dokumenty dzielone są na krótkie fragmenty.
  \item Każdy fragment zamieniany jest na wektor liczb (tak zwane osadzenie),
        który oddaje jego znaczenie. Wektory trafiają do bazy wektorowej.
  \item Pytanie użytkownika również zamieniane jest na wektor.
  \item System wybiera $k$ fragmentów najbardziej podobnych do pytania (zwykle przez
        podobieństwo kosinusowe wektorów).
  \item Te fragmenty trafiają do okna kontekstu modelu, który na ich podstawie formułuje
        odpowiedź.
\end{enumerate}

Całe podejście opiera się więc na jednej operacji: znajdź teksty podobne do pytania.

\subsection{Kiedy to wystarcza}

RAG wektorowy sprawdza się dobrze, gdy odpowiedź mieści się w jednym albo kilku fragmentach
tekstu i gdy pytanie jest podobne do tego fragmentu pod względem słów lub znaczenia. Dobrze
pasują tu proste pytania o fakty, na przykład \enquote{kiedy odbył się szczyt energetyczny
w Warszawie?} albo \enquote{gdzie miała miejsce powódź?}. W takich przypadkach zwykle
istnieje konkretny fragment tekstu, który wprost zawiera odpowiedź, a zadaniem systemu jest
po prostu go odnaleźć.

\subsection{Gdzie zawodzi}

Problem pojawia się przy pytaniach, które wymagają zrozumienia powiązań między elementami
rozrzuconymi po wielu dokumentach. Można je z grubsza podzielić na dwa rodzaje:

\begin{itemize}[itemsep=3pt]
  \item Pytania relacyjne, na przykład \enquote{co łączy Marię Nowak z firmą GreenGrid
        Europe?} albo \enquote{kto brał udział w tym samym wydarzeniu co Amber Analytics?}.
  \item Pytania o sens całości, na przykład \enquote{jakie główne wątki przewijają się
        przez wszystkie te artykuły?} albo \enquote{które zdarzenia łączą się przez tych
        samych uczestników?}.
\end{itemize}

W obu przypadkach nie istnieje pojedynczy fragment, który zawierałby gotową odpowiedź.
Odpowiedź trzeba złożyć z wielu artykułów, często łącząc fakty, które nigdy nie pojawiają się
razem w jednym fragmencie. Tu RAG wektorowy ma fundamentalne ograniczenie. Indeks wektorowy
przechowuje listę fragmentów i potrafi mierzyć ich podobieństwo, ale sama taka struktura nie
jest w stanie naturalnie oddać powiązań między obiektami opisanymi w korpusie. Dwa artykuły,
które wspominają tę samą osobę, są w takim indeksie zupełnie niepowiązane, jeśli używają
różnych słów. System może co najwyżej liczyć, że oba fragmenty
naraz okażą się podobne do pytania i trafią do kontekstu, co przy większej liczbie powiązań
przestaje działać.

\subsection{Pomysł grafu wiedzy}

Systemy grafowe rozwiązują ten problem, zmieniając to, co przechowuje indeks. Zamiast samej
listy fragmentów budują graf wiedzy, w którym:

\begin{itemize}[itemsep=2pt]
  \item węzłami są encje, to znaczy konkretne osoby, organizacje, miejsca i zdarzenia,
  \item krawędziami są relacje między nimi, na przykład \enquote{uczestniczy w},
        \enquote{dostarcza}, \enquote{odbywa się w}.
\end{itemize}

Kluczowa zmiana polega na tym, że relacja staje się pełnoprawnym elementem indeksu, a nie
czymś, co trzeba dopiero odgadnąć z tekstu w trakcie zapytania. Jeśli Maria Nowak pojawia
się w pięciu artykułach, w grafie odpowiada jej jeden węzeł, do którego prowadzą krawędzie ze
wszystkich tych artykułów. Dzięki temu, żeby odpowiedzieć na pytanie o powiązania, system nie
musi mieć szczęścia przy doborze podobnych fragmentów. Może po prostu przejść po krawędziach
grafu i odczytać, co i z czym jest połączone. Rysunek~\ref{fig:pipeline} zestawia oba
podejścia.

% ── Rysunek A: klasyczny RAG vs graph RAG ─────────────────────────────────────
\begin{figure}[h!]
\centering
\begin{tikzpicture}[
  node distance=0.7cm,
  box/.style={draw, rounded corners=3pt, align=center, minimum height=0.85cm,
              minimum width=3.3cm, font=\small},
  cb/.style={box, fill=LightBlue, draw=Blue},
  gb/.style={box, fill=LightTeal, draw=Teal},
  ans/.style={box, font=\small\bfseries},
  arr/.style={-{Stealth[length=2.2mm]}, thick},
  lab/.style={font=\scriptsize\itshape, text=Slate}
]

% lewa kolumna: klasyczny RAG
\node[cb] (d1) {Dokumenty};
\node[cb, below=of d1] (c1) {Fragmenty tekstu};
\node[cb, below=of c1] (v1) {Baza wektorowa};
\node[cb, below=of v1] (t1) {$k$ podobnych fragmentów};
\node[ans, fill=LightBlue, draw=Blue, below=of t1] (a1) {Odpowiedź modelu};

\draw[arr, Blue] (d1) -- node[lab, right]{podział} (c1);
\draw[arr, Blue] (c1) -- node[lab, right]{osadzenia} (v1);
\draw[arr, Blue] (v1) -- node[lab, right]{podobieństwo} (t1);
\draw[arr, Blue] (t1) -- (a1);
\node[font=\bfseries\small, color=Blue, above=0.25cm of d1] {Klasyczny RAG};

% prawa kolumna: graph RAG
\begin{scope}[xshift=7.2cm]
\node[gb] (d2) {Dokumenty};
\node[gb, below=of d2] (c2) {Fragmenty tekstu};
\node[gb, below=of c2] (k2) {Graf wiedzy};
\node[gb, below=of k2] (s2) {Podgraf / społeczności};
\node[ans, fill=LightTeal, draw=Teal, below=of s2] (a2) {Odpowiedź modelu};

\draw[arr, Teal] (d2) -- node[lab, right]{podział} (c2);
\draw[arr, Teal] (c2) -- node[lab, right]{ekstrakcja encji} (k2);
\draw[arr, Teal] (k2) -- node[lab, right]{przejście / podsumowania} (s2);
\draw[arr, Teal] (s2) -- (a2);
\node[font=\bfseries\small, color=Teal, above=0.25cm of d2] {RAG grafowy};
\end{scope}

% zapytanie
\node[box, fill=LightAmber, draw=Amber, minimum width=2.4cm,
      at={($(t1)!0.5!(s2)$)}, yshift=0.0cm] (q) {Pytanie};
\draw[arr, Amber] (q) -- (t1);
\draw[arr, Amber] (q) -- (s2);

\end{tikzpicture}
\caption{Klasyczny RAG wyszukuje fragmenty tekstu podobne do pytania. RAG grafowy buduje z
tych samych dokumentów graf wiedzy i odpowiada, przechodząc po jego węzłach i krawędziach
albo korzystając z gotowych podsumowań całych powiązanych ze sobą grup obiektów.}
\label{fig:pipeline}
\end{figure}

To samo widać dobrze na przykładzie naszego korpusu artykułów prasowych.
Rysunek~\ref{fig:kg} pokazuje uproszczony wycinek grafu wiedzy zbudowanego z wiadomości o
bałtyckiej sieci energetycznej i o powodzi w Krakowie (korpus opisany jest dokładniej w
dodatku~\ref{app:korpus}). Widać tam powiązania, których żaden
pojedynczy artykuł nie zawiera w całości. Na przykład firma Amber Analytics dostarcza
czujniki do projektu Baltic Grid Agreement, a jednocześnie drony użyte przy akcji
powodziowej w Krakowie. Podobnie Maria Nowak pojawia się i na szczycie energetycznym, i przy
reakcji na powódź. To są dokładnie te mosty między wątkami, które wyszukiwanie po
podobieństwie tekstu gubi, a które w grafie są po prostu krawędziami.

% ── Rysunek B: graf wiedzy z korpusu ──────────────────────────────────────────
\begin{figure}[h!]
\centering
\begin{tikzpicture}[
  every node/.style={font=\footnotesize},
  ent/.style={draw, rounded corners=3pt, align=center, inner sep=4pt, font=\footnotesize},
  person/.style={ent, fill=LightBlue, draw=Blue},
  org/.style={ent, fill=LightTeal, draw=Teal},
  event/.style={ent, fill=LightAmber, draw=Amber},
  loc/.style={ent, fill=LightSlate, draw=Slate},
  rel/.style={-{Stealth[length=1.8mm]}, draw=Slate!80, thin},
  bridge/.style={-{Stealth[length=1.8mm]}, draw=Amber, thick, dashed},
  rl/.style={font=\scriptsize\itshape, text=Slate, fill=white, inner sep=1.5pt}
]

% wątek energetyczny (lewa strona)
\node[event]  (summit) at (1.7, 6.6) {Szczyt energetyczny\\w Warszawie};
\node[loc]    (wawa)   at (0.4, 4.9) {Warszawa};
\node[person] (nowak)  at (6.4, 7.1) {Maria Nowak};
\node[event]  (bga)    at (3.4, 4.3) {Baltic Grid\\Agreement};
\node[org]    (green)  at (1.4, 2.6) {GreenGrid Europe};
\node[org]    (amber)  at (4.9, 2.4) {Amber Analytics};

% wątek powodziowy (prawa strona)
\node[event]  (flood)  at (9.3, 4.1) {Akcja powodziowa\\w Krakowie};
\node[loc]    (krakow) at (10.9, 2.4) {Kraków};

% relacje w obrębie wątku
\draw[rel] (nowak)  -- node[rl, pos=0.45]{uczestniczy} (summit);
\draw[rel] (nowak)  -- node[rl]{ogłasza} (bga);
\draw[rel] (green)  -- node[rl]{realizuje} (bga);
\draw[rel] (amber)  -- node[rl]{czujniki} (bga);
\draw[rel] (summit) -- node[rl]{odbywa się w} (wawa);
\draw[rel] (flood)  -- node[rl]{miejsce} (krakow);

% mosty między wątkami
\draw[bridge] (nowak) -- node[rl, pos=0.62]{komentuje} (flood);
\draw[bridge] (amber) -- node[rl, pos=0.5]{drony} (flood);

% legenda
\begin{scope}[shift={(0,0.3)}]
  \node[person, minimum width=1.4cm] (lp) at (0.6,0) {osoba};
  \node[org, minimum width=1.4cm, right=0.2cm of lp] (lo) {organizacja};
  \node[event, minimum width=1.4cm, right=0.2cm of lo] (le) {zdarzenie};
  \node[loc, minimum width=1.4cm, right=0.2cm of le] (ll) {lokalizacja};
  \draw[bridge] ([xshift=0.3cm]ll.east) -- ++(0.8,0)
        node[right, font=\scriptsize, text=Slate]{powiązanie między wątkami};
\end{scope}

\end{tikzpicture}
\caption{Uproszczony przykład grafu wiedzy zbudowanego z naszego korpusu artykułów. Węzły to
obiekty różnych typów, krawędzie to nazwane relacje. Linie przerywane wyróżniają powiązania,
które łączą oddzielne wątki (energetyczny i powodziowy) i których nie znajdziemy w żadnym
pojedynczym artykule.}
\label{fig:kg}
\end{figure}

% ══════════════════════════════════════════════════════════════════════════════
\section{GraphRAG firmy Microsoft}
% ══════════════════════════════════════════════════════════════════════════════

GraphRAG opisany przez Edge i współpracowników \cite{edge2024} powstał, by rozwiązać
konkretny problem, z którym klasyczny RAG sobie nie radzi: globalne rozumienie całego zbioru
dokumentów. Chodzi o pytania w stylu \enquote{jakie są główne wątki w tych materiałach?}, na
które nie ma jednego najlepszego fragmentu, bo odpowiedź wymaga streszczenia całości. Autorzy
nazywają to zadaniem podsumowania ukierunkowanego pytaniem.

System działa w dwóch fazach. Najpierw, raz, buduje indeks z dokumentów. Potem korzysta z
tego indeksu przy każdym zapytaniu. Warto je rozdzielić, bo to właśnie kosztowna faza
indeksowania odróżnia GraphRAG od prostszego RAG.

\subsection{Faza 1: indeksowanie}

Budowa indeksu przebiega w kilku krokach, od surowego tekstu do gotowych podsumowań.

\subsubsection{Podział na fragmenty}

Dokumenty dzielone są na zachodzące na siebie fragmenty (w pracy domyślnie 600 tokenów ze
100-tokenowym zakładem). Rozmiar fragmentu to istotny wybór projektowy. Dłuższe fragmenty
oznaczają mniej wywołań modelu podczas ekstrakcji, ale model gorzej wychwytuje informacje
pojawiające się pod koniec długiego fragmentu.

\subsubsection{Ekstrakcja encji i relacji}

Dla każdego fragmentu model językowy jest proszony o wypisanie:

\begin{itemize}[itemsep=2pt]
  \item encji, czyli nazwanych osób, organizacji, miejsc i zdarzeń, razem z krótkim opisem,
  \item relacji, czyli skierowanych powiązań między parami encji, każdej z opisem,
  \item faktów, czyli krótkich, weryfikowalnych stwierdzeń o encjach, na przykład dat i ról.
\end{itemize}

Prompt można dopasować do dziedziny, dodając przykłady charakterystyczne dla nauki, prawa
czy medycyny. Ten krok jest sercem indeksowania i zarazem jego głównym kosztem, bo model
trzeba wywołać osobno dla każdego fragmentu.

\subsubsection{Budowa grafu wiedzy}

Wszystkie wydobyte encje i relacje łączone są w jeden skierowany graf. Ta sama encja często
pojawia się w wielu fragmentach, więc jej wystąpienia są scalane, a opisy podsumowywane przez
model. Powtarzające się relacje również są łączone, a waga krawędzi odpowiada temu, jak
często dana relacja występowała w tekście. W podstawowej wersji dopasowanie nazw odbywa się
przez dokładne porównanie napisów, choć możliwe jest też łagodniejsze dopasowanie.

\subsubsection{Wykrywanie społeczności}

To krok, który najmocniej odróżnia GraphRAG od innych podejść grafowych. Na gotowym grafie
uruchamiany jest algorytm wykrywania społeczności \textit{Leiden} \cite{traag2019}. Dzieli on
graf na grupy gęsto powiązanych encji, na przykład wszystkie wątki wokół jednego szczytu albo
wszystkich uczestników jednej afery. Co ważne, robi to hierarchicznie i na wielu poziomach:

\begin{itemize}[itemsep=2pt]
  \item poziom C0 to najszersze, najbardziej ogólne społeczności (jest ich najmniej),
  \item poziomy C1, C2, C3 to coraz drobniejsze podspołeczności.
\end{itemize}

Każdy poziom pokrywa cały graf w sposób rozłączny i kompletny, więc daje inną
\enquote{rozdzielczość} spojrzenia na zbiór. Rysunek~\ref{fig:comm} pokazuje tę ideę.

\subsubsection{Podsumowania społeczności}

Dla każdej społeczności, na każdym poziomie, model generuje opisowe podsumowanie w języku
naturalnym. Podsumowania najniższego poziomu powstają z opisów encji, relacji i faktów.
Podsumowania wyższych poziomów powstają z podsumowań ich podspołeczności, oddolnie. W efekcie
powstaje bogaty, hierarchiczny zbiór gotowych streszczeń, które obejmują cały korpus jeszcze
zanim padnie jakiekolwiek pytanie. To te streszczenia pozwalają później odpowiadać na pytania
o całość bez czytania wszystkich fragmentów na nowo.

% ── Rysunek C: hierarchia społeczności ────────────────────────────────────────
\begin{figure}[h!]
\centering
\begin{tikzpicture}[
  every node/.style={font=\footnotesize},
  dot/.style={circle, draw, minimum size=4.5pt, inner sep=0pt},
  tnode/.style={draw, rounded corners=2pt, align=center, font=\scriptsize,
                minimum height=0.55cm, minimum width=1.5cm},
  arr/.style={-{Stealth[length=1.6mm]}, draw=Slate}
]

% lewa strona: graf z trzema społecznościami
\begin{scope}
  % społeczność 1 (teal) - sieć energetyczna
  \node[dot, fill=Teal] (g1) at (0.3,3.2) {};
  \node[dot, fill=Teal] (g2) at (1.1,3.7) {};
  \node[dot, fill=Teal] (g3) at (1.6,2.9) {};
  \node[dot, fill=Teal] (g4) at (0.8,2.6) {};
  \draw[Teal] (g1)--(g2) (g2)--(g3) (g3)--(g4) (g4)--(g1) (g1)--(g3);

  % społeczność 2 (amber) - reagowanie kryzysowe
  \node[dot, fill=Amber] (a1) at (3.2,1.0) {};
  \node[dot, fill=Amber] (a2) at (3.9,1.5) {};
  \node[dot, fill=Amber] (a3) at (4.3,0.7) {};
  \draw[Amber] (a1)--(a2) (a2)--(a3) (a3)--(a1);

  % społeczność 3 (blue) - administracja
  \node[dot, fill=Blue] (b1) at (3.6,3.6) {};
  \node[dot, fill=Blue] (b2) at (4.4,3.9) {};
  \node[dot, fill=Blue] (b3) at (4.2,3.0) {};
  \draw[Blue] (b1)--(b2) (b2)--(b3) (b3)--(b1);

  % krawędzie między społecznościami
  \draw[Slate!50] (g3)--(b3) (a2)--(b3);

  \begin{pgfonlayer}{background}
    \node[fill=LightTeal, rounded corners=6pt, fit=(g1)(g2)(g3)(g4), inner sep=5pt] {};
    \node[fill=LightAmber, rounded corners=6pt, fit=(a1)(a2)(a3), inner sep=5pt] {};
    \node[fill=LightBlue, rounded corners=6pt, fit=(b1)(b2)(b3), inner sep=5pt] {};
  \end{pgfonlayer}
  \node[font=\scriptsize\itshape, text=Slate] at (2.3,4.5) {graf podzielony na społeczności};
\end{scope}

% prawa strona: hierarchia + podsumowania
\begin{scope}[xshift=7.0cm, yshift=0.2cm]
  \node[tnode, fill=LightSlate] (c0) at (1.5,4.2) {C0: cały korpus};
  \node[tnode, fill=LightTeal] (c1a) at (0,2.6) {C1: energia};
  \node[tnode, fill=LightBlue] (c1b) at (3,2.6) {C1: kryzys};
  \node[tnode, fill=LightTeal] (c2a) at (-0.6,1.0) {C2};
  \node[tnode, fill=LightTeal] (c2b) at (1.0,1.0) {C2};
  \node[tnode, fill=LightBlue] (c2c) at (3,1.0) {C2};

  \draw[arr] (c0)--(c1a);
  \draw[arr] (c0)--(c1b);
  \draw[arr] (c1a)--(c2a);
  \draw[arr] (c1a)--(c2b);
  \draw[arr] (c1b)--(c2c);
  \node[font=\scriptsize\itshape, text=Slate, align=center] at (1.5,5.0)
    {hierarchia społeczności\\(każdy węzeł ma podsumowanie od modelu)};
\end{scope}

\end{tikzpicture}
\caption{Wykrywanie społeczności dzieli graf na grupy gęsto powiązanych encji (po lewej), a
robi to na wielu poziomach jednocześnie (po prawej). Dla każdej społeczności na każdym
poziomie model generuje osobne podsumowanie. Poziom C0 daje najszersze spojrzenie, niższe
poziomy coraz drobniejsze.}
\label{fig:comm}
\end{figure}

\subsubsection{Osadzenia wektorowe}

Warto w tym miejscu rozwiać częste nieporozumienie. GraphRAG nie rezygnuje z klasycznych
osadzeń wektorowych, tych samych, na których opiera się zwykły RAG. Oprócz grafu i podsumowań
system liczy na etapie indeksowania osadzenia dla opisów encji oraz dla podsumowań
społeczności i zapisuje je w bazie wektorowej. Pełnią one jednak rolę pomocniczą, a nie
główną. Służą do tego, żeby przy zapytaniu szybko dopasować pytanie do właściwych encji albo
podsumowań, czyli znaleźć dobry punkt startowy w grafie. Sama odpowiedź powstaje już z treści
grafu i streszczeń, a nie z surowych fragmentów wybranych po podobieństwie.
Z tych osadzeń korzysta wyszukiwanie lokalne oraz DRIFT.

\subsection{Faza 2: tryby zapytań}

GraphRAG udostępnia kilka trybów odpowiadania, które korzystają z gotowego indeksu na zupełnie
różne sposoby. Tryb wybiera się dla każdego pytania osobno.

\subsubsection{Wyszukiwanie lokalne}

Wyszukiwanie lokalne służy do pytań, które wskazują konkretną encję,
na przykład \enquote{co zrobiła firma GreenGrid Europe?}. System zamienia pytanie na wektor i
porównuje je z policzonymi wcześniej osadzeniami opisów encji, żeby znaleźć węzły startowe.
Jest to więc dokładnie ten sam mechanizm, co w zwykłym RAG, tyle że tu szuka się pasujących
encji, a nie pasujących fragmentów tekstu. Następnie system zbiera wokół
nich kontekst z kilku źródeł naraz: fragmenty tekstu wspominające te encje, podsumowania
społeczności, do których encje należą, sąsiednie encje w grafie, krawędzie wychodzące z
węzłów startowych oraz związane z nimi fakty. To wszystko jest porządkowane według trafności,
upakowywane w jedno okno kontekstu i przekazywane do modelu, który jednym wywołaniem tworzy
odpowiedź. Lokalne wyszukiwanie jest więc tanie, bo ogranicza się do najbliższego otoczenia
wskazanych encji.

\subsubsection{Wyszukiwanie globalne}

Wyszukiwanie globalne odpowiada na pytania bez konkretnej kotwicy,
które dotyczą całego korpusu. Co istotne, ten tryb w ogóle nie sięga do grafu w trakcie
zapytania i nie korzysta z osadzeń wektorowych. Działa wyłącznie na wcześniej przygotowanych
podsumowaniach społeczności, w schemacie map-reduce:

\begin{enumerate}[itemsep=2pt]
  \item Wybierany jest poziom społeczności (na przykład C0 albo C2) i pobierane są wszystkie
        podsumowania z tego poziomu.
  \item Podsumowania są tasowane (żeby uniknąć faworyzowania początku kontekstu) i dzielone
        na porcje mieszczące się w oknie modelu.
  \item Krok \textit{map}: każda porcja trafia osobno do modelu, który wypisuje listę
        krótkich stwierdzeń istotnych dla pytania, każde z oceną ważności od 0 do 100.
        Wszystkie porcje przetwarzane są równolegle.
  \item Stwierdzenia z oceną 0 są odrzucane.
  \item Krok \textit{reduce}: pozostałe stwierdzenia są sortowane według ważności, pakowane
        w nowe okno kontekstu i jednym końcowym wywołaniem model składa z nich globalną
        odpowiedź.
\end{enumerate}

To podejście jest skuteczne, ale kosztowne, bo w kroku \textit{map} liczba wywołań modelu
rośnie wraz z liczbą społeczności na wybranym poziomie. Na największym zbiorze z pracy
poziom drobny oznacza setki wywołań na jedno pytanie, podczas gdy poziom C0 wystarcza zaledwie
kilkudziesięciu wywołaniom przy umiarkowanym spadku jakości.

\subsubsection{Wyszukiwanie DRIFT}

Istnieje też tryb pośredni, DRIFT, który zaczyna od kontekstu globalnego (podsumowań
społeczności), a potem iteracyjnie schodzi w głąb grafu poprzez wyszukiwanie lokalne,
zadając kolejne pytania pomocnicze i pogłębiając te wątki, w których jest najbardziej pewny.
Łączy więc szerokie spojrzenie z lokalną dokładnością.

\subsection{Wyniki z pracy}

Autorzy porównali GraphRAG z klasycznym RAG wektorowym na dwóch dużych korpusach (transkrypcje
podcastów i artykuły prasowe, każdy rzędu miliona tokenów), używając modelu jako sędziego
oceniającego odpowiedzi. GraphRAG na wszystkich poziomach społeczności wyraźnie wygrywał pod
względem kompletności i różnorodności odpowiedzi, z przewagą rzędu 72--83\% i 62--82\%. Co
ważne, podsumowania na poziomie C0 osiągały ten wynik, zużywając około 97\% mniej tokenów niż
streszczanie wprost z tekstu źródłowego.

\subsection{Ograniczenia}

GraphRAG płaci za swoją skuteczność wysoką ceną, którą warto wyraźnie nazwać, bo to ona
motywuje powstanie lżejszych systemów takich jak LightRAG.

\begin{itemize}[itemsep=3pt]
  \item \textbf{Kosztowne indeksowanie.} Model trzeba wywołać dla każdego fragmentu przy
        ekstrakcji, a potem jeszcze raz dla każdej społeczności na każdym poziomie przy
        generowaniu podsumowań. To przekłada się na bardzo dużo wywołań i tokenów, a w pracy
        zbudowanie indeksu zajmowało rzędu kilku godzin.
  \item \textbf{Drogie zapytania globalne.} Schemat map-reduce wykonuje osobne wywołanie
        modelu na społeczność, więc na większych zbiorach jedno pytanie globalne to setki
        wywołań.
  \item \textbf{Trudna aktualizacja.} To może być najpoważniejsze ograniczenie w praktyce.
        Dodanie nowych artykułów zmienia strukturę grafu, więc trzeba na nowo wykryć
        społeczności i wygenerować dotknięte podsumowania. Przy strumieniu napływających
        wiadomości oznacza to ciągłe, kosztowne przebudowywanie indeksu.
\end{itemize}

% ══════════════════════════════════════════════════════════════════════════════
\section{LightRAG}
% ══════════════════════════════════════════════════════════════════════════════

LightRAG \cite{guo2024} to system zaprojektowany wprost po to, by usunąć dwa największe
koszty GraphRAG opisane wyżej: drogie zapytania i kosztowną aktualizację indeksu.
Najkrócej: LightRAG rezygnuje z hierarchii społeczności i z generowania ich
podsumowań, a w zamian wprowadza dwupoziomowe wyszukiwanie po grafie wzbogaconym o bazę
wektorową.

\subsection{Indeksowanie}

LightRAG buduje graf wiedzy podobnie jak GraphRAG, w trzech operacjach na każdym fragmencie:

\begin{itemize}[itemsep=3pt]
  \item \textbf{Ekstrakcja.} Model wydobywa encje i relacje z fragmentu, tak jak w GraphRAG.
  \item \textbf{Profilowanie do par klucz-wartość.} Tu pojawia się główna różnica. Każdy
        węzeł i każda krawędź zamieniane są na parę $(K, V)$. Klucz $K$ to słowo lub krótka
        fraza służąca do wyszukiwania, a wartość $V$ to wygenerowany przez model akapit
        opisujący daną encję lub relację wraz z istotnymi fragmentami źródła. Encja ma jeden
        klucz, którym jest jej nazwa. Relacja może mieć kilka kluczy, bo są nimi
        ogólniejsze frazy tematyczne opisujące, czego dotyczy powiązanie, na przykład
        \enquote{współpraca energetyczna} czy \enquote{nadzór regulacyjny}.
  \item \textbf{Deduplikacja.} Identyczne encje i relacje z różnych fragmentów są scalane,
        żeby graf pozostał zwarty.
\end{itemize}

Najważniejsze, czego tu nie ma: nie ma wykrywania społeczności ani generowania ich
podsumowań. To właśnie ten krok był w GraphRAG najdroższy i najtrudniejszy w aktualizacji.
Klucze $K$ zapisywane są jako wektory w bazie wektorowej, więc do encji i relacji można potem
docierać przez podobieństwo, podobnie jak do fragmentów w klasycznym RAG, tyle że tu obiektem
wyszukiwania jest struktura grafu, a nie surowy tekst.

\subsubsection{Aktualizacja przyrostowa}

Nowy dokument przechodzi przez te same trzy operacje, a powstałe węzły i krawędzie są po
prostu dołączane do istniejącego grafu jako suma zbiorów. Nie trzeba przebudowywać żadnej
struktury społeczności, bo jej nie ma. To bezpośrednie rozwiązanie problemu, który w GraphRAG
wymagał kosztownego przeliczania podsumowań przy każdej zmianie danych. Dla zbioru artykułów,
który stale rośnie, jest to istotna różnica praktyczna.

\subsection{Dwupoziomowe wyszukiwanie}

Pomysł polega na tym, że pytania bywają dwojakiego rodzaju, więc i wyszukiwanie ma dwa
poziomy:

\begin{itemize}[itemsep=2pt]
  \item poziom niski dotyczy konkretnych encji i ich atrybutów, na przykład
        \enquote{kim jest Maria Nowak?},
  \item poziom wysoki dotyczy szerszych tematów i powiązań, na przykład
        \enquote{jak współpraca energetyczna wpływa na region?}.
\end{itemize}

Samo wyszukiwanie wygląda następująco:

\begin{enumerate}[itemsep=3pt]
  \item \textbf{Wydobycie słów kluczowych (jedno wywołanie modelu).} Pytanie trafia do
        modelu, który wypisuje dwie listy: słowa lokalne (nazwy konkretnych encji) oraz słowa
        globalne (frazy tematyczne opisujące, o co w pytaniu chodzi).
  \item \textbf{Dopasowanie niskiego poziomu.} Słowa lokalne są porównywane z kluczami encji
        w bazie wektorowej, co zwraca wartości $V$ pasujących encji.
  \item \textbf{Dopasowanie wysokiego poziomu.} Słowa globalne są porównywane z kluczami
        tematycznymi relacji, co zwraca wartości $V$ pasujących powiązań.
  \item \textbf{Rozszerzenie o sąsiadów.} Dla dopasowanych encji i relacji dołączani są ich
        bezpośredni sąsiedzi w grafie (sąsiedzi w odległości jednej krawędzi). Dzięki temu w
        kontekście pojawiają się też encje, których nazwa nie padła w pytaniu, ale które są
        powiązane z tymi, które padły.
  \item \textbf{Złożenie kontekstu i odpowiedź (jedno wywołanie modelu).} Wszystkie zebrane
        wartości $V$ są łączone i przekazywane modelowi razem z pytaniem, a model jednym
        wywołaniem tworzy odpowiedź.
\end{enumerate}

W trybie hybrydowym, który jest zalecanym domyślnym ustawieniem, kroki dopasowania niskiego i
wysokiego poziomu działają jednocześnie, więc system obsługuje zarówno pytania o konkretne
encje, jak i o ogólne tematy. Kluczowa obserwacja kosztowa jest taka, że całe wyszukiwanie
mieści się w zaledwie dwóch wywołaniach modelu, niezależnie od wielkości korpusu: jedno na
wydobycie słów kluczowych i jedno na odpowiedź. Nie ma tu przechodzenia po setkach
społeczności, bo społeczności w ogóle nie istnieją.

Ciekawym wynikiem pracy jest test, w którym z wartości $V$ usunięto oryginalny
tekst źródłowy, zostawiając tylko streszczenia wygenerowane przez model. Jakość
odpowiedzi nie spadła, a w części dziedzin nawet wzrosła. Oznacza to, że
uporządkowane streszczenia encji i relacji potrfią nieść wystarczająco dużo
informacji by zastąpić surowy tekst.

\subsection{Porównanie z GraphRAG}

LightRAG i GraphRAG dzielą tę samą podstawową intuicję, że strukturę powiązań warto zapisać
wprost przy budowie indeksu. Różnią się tym, co robią dalej. GraphRAG dokłada hierarchię
społeczności i ich podsumowania, co świetnie służy pytaniom o sens całości, ale jest drogie i
trudne w aktualizacji. LightRAG rezygnuje ze społeczności na rzecz dwupoziomowego
wyszukiwania po parach klucz-wartość, co daje stały, niski koszt zapytania i tanie
aktualizacje. Tabela~\ref{tab:porownanie} zbiera te różnice.

\begin{table}[h!]
\centering
\renewcommand{\arraystretch}{1.35}
\begin{tabularx}{\textwidth}{@{} p{3.4cm} X X @{}}
\toprule
\textbf{Aspekt} & \textbf{GraphRAG (Microsoft)} & \textbf{LightRAG} \\
\midrule
Główny pomysł
  & hierarchia społeczności i podsumowania, odpowiedź metodą map-reduce
  & płaski graf encji i relacji, odpowiedź przez dwupoziomowe wyszukiwanie \\
Struktura grafu
  & graf wiedzy z wielopoziomową hierarchią społeczności
  & graf wiedzy z bazą wektorową par klucz-wartość \\
Wywołania modelu przy zapytaniu
  & rośnie z liczbą społeczności (setki dla dużych zbiorów)
  & stałe, dwa wywołania niezależnie od wielkości korpusu \\
Aktualizacja danych
  & wymaga ponownego wykrycia społeczności i podsumowań
  & dołączenie nowych węzłów i krawędzi, prawie bez kosztu \\
Najlepszy typ pytań
  & globalne pytania o sens całości
  & zarówno konkretne, jak i tematyczne pytania o powiązania \\
Słaby punkt
  & wysoki koszt i trudna aktualizacja
  & słabszy w czysto globalnym streszczaniu niż podsumowania społeczności \\
\bottomrule
\end{tabularx}
\caption{Zestawienie GraphRAG i LightRAG na najważniejszych wymiarach.}
\label{tab:porownanie}
\end{table}

W eksperymentach na czterech dziedzinach LightRAG wypadał lepiej od wszystkich porównywanych
podejść, w tym od GraphRAG, jednocześnie będąc znacznie tańszym przy zapytaniach. Różnica w
koszcie jest wyraźna: na największym zbiorze jedno zapytanie w GraphRAG zużywało setki
tysięcy tokenów i setki wywołań modelu, podczas gdy w LightRAG było to mniej niż sto tokenów
i pojedyncze wywołanie wyszukiwania. To pokazuje, że rezygnacja ze społeczności nie musi
oznaczać gorszych odpowiedzi, a daje duże oszczędności tam, gdzie dane często się zmieniają.

% ══════════════════════════════════════════════════════════════════════════════
\section{Podsumowanie}
% ══════════════════════════════════════════════════════════════════════════════

Klasyczny RAG opiera się na jednej operacji, czyli na wyszukiwaniu tekstów podobnych do
pytania. Dla prostych pytań o fakty to w zupełności wystarcza i nie ma powodu sięgać po nic
więcej. Trudność zaczyna się przy pytaniach o powiązania i o sens całości, bo taki indeks nie
oddaje wprost powiązań między obiektami i potrafi je odtworzyć tylko przypadkiem.

Systemy GraphRAG rozwiązują to, zapisując encje i relacje wprost w grafie wiedzy już na etapie
indeksowania. GraphRAG firmy Microsoft idzie krok dalej i buduje hierarchię społeczności wraz
z gotowymi podsumowaniami, dzięki czemu potrafi odpowiadać na pytania o całość korpusu. Płaci
za to wysokim kosztem indeksowania, drogimi zapytaniami globalnymi i trudną aktualizacją
przy nowych danych. LightRAG zachowuje sam graf, ale rezygnuje ze społeczności na rzecz
dwupoziomowego wyszukiwania po parach klucz-wartość. Dzięki temu utrzymuje stały, niski koszt
zapytania i pozwala tanio dokładać nowe dokumenty, co jest istotne dla stale rosnącego zbioru
artykułów prasowych.

Wybór między tymi podejściami zależy więc od zastosowania. Tam, gdzie najważniejsze są
globalne streszczenia dużego, w miarę stałego korpusu, naturalnie pasuje GraphRAG. Tam, gdzie
dane często się zmieniają, a pytania mieszają konkretne encje z ogólnymi tematami, lepiej
sprawdza się LightRAG. Klasyczny RAG pozostaje rozsądnym, tanim punktem odniesienia dla
prostych pytań. Te obserwacje stanowią tło dla naszych własnych eksperymentów na korpusie
artykułów, które opiszemy w kolejnej części dokumentacji.

% ══════════════════════════════════════════════════════════════════════════════
\appendix
\section{Korpus użyty w projekcie}
\label{app:korpus}
% ══════════════════════════════════════════════════════════════════════════════

Do testów i przykładów w tym raporcie wykorzystujemy niewielki, ale spójny korpus wiadomości
prasowych. Składa się on z 18 artykułów, oznaczonych jako materiał syntetyczny i opatrzonych
datami z okresu od kwietnia do maja 2026 roku. Artykuły zostały przygotowane tak, aby tworzyły
jedną, powiązaną historię, a nie zbiór niezależnych tekstów.

Korpus splata dwa wątki:
\begin{itemize}[itemsep=2pt]
  \item wątek energetyczny wokół projektu rozbudowy bałtyckiej sieci energetycznej (Baltic
        Grid Agreement), szczytu w Warszawie, morskich farm wiatrowych pod Gdańskiem oraz
        magazynów energii,
  \item wątek kryzysowy wokół powodzi w Krakowie i akcji ratunkowej prowadzonej przez
        krajowe służby.
\end{itemize}

Najważniejszą cechą korpusu jest to, że te same obiekty powracają w wielu artykułach i łączą
oba wątki. Na przykład firma Amber Analytics dostarcza czujniki do projektu energetycznego, a
zarazem drony wykorzystane podczas powodzi, zaś Maria Nowak występuje i na szczycie
energetycznym, i przy reakcji na kryzys. Dzięki temu korpus zawiera realne powiązania
wielokrokowe, których nie da się odczytać z pojedynczego artykułu, co czyni go dobrym
materiałem do porównania wyszukiwania klasycznego z grafowym.

W korpusie występują obiekty czterech głównych typów:
\begin{itemize}[itemsep=2pt]
  \item osoby, na przykład Maria Nowak, Jonas Berg, Ewa Lewandowska,
  \item organizacje, na przykład GreenGrid Europe, NordicGrid, Amber Analytics, Clear Air
        Forum,
  \item zdarzenia, na przykład Baltic Grid Agreement, szczyt energetyczny w Warszawie, powódź
        w Krakowie,
  \item lokalizacje, na przykład Warszawa, Gdańsk, Wilno, Kraków.
\end{itemize}

Do korpusu dołączony jest zestaw 18 pytań w stylu śledczym, na przykład \enquote{co łączy
Marię Nowak z firmą GreenGrid Europe?}, z których każde ma przypisaną oczekiwaną odpowiedź.
Służą one do oceny poszczególnych trybów wyszukiwania. 

% ── bibliografia ──────────────────────────────────────────────────────────────
\begin{thebibliography}{9}

\bibitem{edge2024}
Edge, D., Trinh, H., Cheng, N., Bradley, J., Chao, A., Mody, A., Truitt, S.,
Metropolitansky, D., Ness, R.~O., Larson, J. (2024).
\textit{From Local to Global: A GraphRAG Approach to Query-Focused Summarization.}
arXiv:2404.16130.

\bibitem{guo2024}
Guo, Z., Xia, L., Yu, Y., Ao, T., Huang, C. (2024).
\textit{LightRAG: Simple and Fast Retrieval-Augmented Generation.}
arXiv:2410.05779.

\bibitem{traag2019}
Traag, V.~A., Waltman, L., van Eck, N.~J. (2019).
\textit{From Louvain to Leiden: guaranteeing well-connected communities.}
Scientific Reports, 9, 5233.

\end{thebibliography}

\end{document}
